\documentclass[11pt]{article}
\usepackage[margin=1.1in]{geometry}
\usepackage{amsmath,amssymb,amsthm,booktabs,hyperref}
\hypersetup{colorlinks=true,linkcolor=blue,citecolor=blue,urlcolor=blue}

\newtheorem{theorem}{Theorem}
\newtheorem{lemma}[theorem]{Lemma}
\newtheorem{proposition}[theorem]{Proposition}
\newtheorem{corollary}[theorem]{Corollary}
\newtheorem{conjecture}[theorem]{Conjecture}
\theoremstyle{definition}
\newtheorem{definition}[theorem]{Definition}
\theoremstyle{remark}
\newtheorem{remark}[theorem]{Remark}

\newcommand{\Q}{\mathbb{Q}}
\newcommand{\La}{\Lambda}
\newcommand{\Pl}{\mathrm{Pl}}
\newcommand{\Res}{\operatorname{Res}}
\newcommand{\Hom}{\operatorname{Hom}}
\newcommand{\Gr}{\operatorname{Gr}}
\newcommand{\soc}{\operatorname{soc}}
\newcommand{\qbin}[2]{\begin{bmatrix}#1\\#2\end{bmatrix}_q}
\newcommand{\fE}{\mathfrak{E}}

\title{Dyck Paths in the Nullspace\\[0.4em]\large A proof of the Clifford--Stanley--Alexandersson conjecture on the
power-sum to monomial transition matrix}
\author{Prajit Chennupati\\ \small\texttt{chennupati@gatech.edu}}
\date{November 27, 2026}

\begin{document}

\maketitle
\begin{abstract}
Let $R_n$ be the matrix expressing the power sums $p_\lambda$ ($\lambda\vdash n$) in the monomial basis, and let
$D_n$ be its diagonal. Clifford and Stanley computed the nullity of $R_n-D_n$ for small $n$ but did not conjecture
its value. Alexandersson conjectured in 2019 that it equals the number of Dyck paths whose peak abscissae sum
to $n$ (OEIS A129528). We prove this conjecture. More precisely, we show that the nullspace, filtered by the
length of partitions, has graded pieces of dimension $[q^{n-d}]\,C_d(q)$, where $C_d(q)$ is MacMahon's
$q$-Catalan number. The proof turns the nullspace into the invariants of a one-parameter group acting on power series.
For the leading part of this group, writing a series as a ratio $u=v/w$ identifies the invariants with
$\Hom_{S_{2d}}(S^{(d,d)},R_{2d})^*$, where $R_{2d}$ is the coinvariant algebra. To control the full group we
prove two further results. First, the top invariant lifts, through an explicit sum over rooted forests. Second,
the whole space of invariants is generated from the top one by reparametrizations of the series variable. The
second result is proved in the cohomology of the Grassmannian $\Gr(d,2d)$, using a screening operator on the
$c=1$ Fock space.
\end{abstract}

\section{The conjecture}

Let $\La$ be the ring of symmetric functions over $\Q$, with power sums $p_\lambda$, monomial symmetric functions
$m_\mu$ and complete homogeneous functions $h_\lambda$ \cite{Mac}. For $\lambda,\mu\vdash n$ define $R_{\lambda\mu}$ by
\[
  p_\lambda=\sum_{\mu\vdash n}R_{\lambda\mu}\,m_\mu ,
\]
and let $D_n=\operatorname{diag}(R_{\lambda\lambda})$. Clifford and Stanley \cite{CS04} met the matrix $R_n-D_n$ in their study of
bottom Schur functions. They computed its nullity for small $n$, obtaining
$1,\allowbreak1,\allowbreak1,\allowbreak2,\allowbreak2,\allowbreak3,\allowbreak4,\allowbreak5,\allowbreak7,\allowbreak9,\allowbreak11,\allowbreak15,\allowbreak19,\allowbreak24,\allowbreak\dots$, and remarked that they had no conjecture for its value. In 2019
Alexandersson observed that these numbers match OEIS A129528 \cite{OEIS}, and recorded the following on the
symmetric functions catalog \cite{SymCat}.

\begin{conjecture}[{Alexandersson, 2019 \cite{SymCat}}]\label{conj}
For every $n\ge1$, $\dim_\Q\ker(R_n-D_n)$ equals the number $a(n)$ of Dyck paths (of any semilength) whose peak
abscissae sum to $n$.
\end{conjecture}

Here a Dyck path of semilength $d$ is a word in $U=(1,1)$ and $D=(1,-1)$ that starts and ends at height $0$ and never
goes below it. A peak is a factor $UD$, and its abscissa is the $x$-coordinate of its top. For example $UD$ has
peak sum $1$, and $UUDD$ has peak sum $2$.

\begin{theorem}\label{main}
Conjecture~\ref{conj} is true: $\dim_\Q\ker(R_n-D_n)=a(n)$ for all $n\ge1$.
\end{theorem}

We prove a refinement. Write $\ell(\lambda)$ for the length of a partition and $C_d(q)=\qbin{2d}{d}/[d+1]_q$
for MacMahon's $q$-Catalan number.

\begin{theorem}\label{refined}
Let $F_{\ge d}\subseteq\La^n$ be the span of the $h_\lambda$ with $\ell(\lambda)\ge d$. Under the identification
of Section~\ref{sec:reduction}, the nullspace of $R_n-D_n$ becomes a subspace $N\subseteq\La^n$ such that
\[
  \dim\bigl(N\cap F_{\ge d}\bigr)/\bigl(N\cap F_{\ge d+1}\bigr)=[q^{\,n-d}]\,C_d(q)\qquad(1\le d\le n).
\]
\end{theorem}

Summing over $d$ gives Theorem~\ref{main}, because $\sum_d q^dC_d(q)$ counts Dyck paths by peak sum
(Section~\ref{sec:dyck}).

\paragraph{Outline.}
\begin{itemize}
\item Section~\ref{sec:reduction} identifies the nullspace with the kernel of a derivation $D_v$ of
$\La=\Q[h_1,h_2,\dots]$. Here $D_v$ is induced by a formal vector field $v(u)\partial_u$ acting on power
series $u=\sum h_kt^k$, and $v$ is the iterative logarithm of $\log(1+u)$.
\item Section~\ref{sec:thmC} computes the kernel of the leading part $u^2\partial_u$ (Theorem~\ref{thmC}),
using Pl\"ucker coordinates, Schur--Weyl duality and Chevalley's theorem.
\item Section~\ref{sec:dyck} gives the Dyck path count.
\item Section~\ref{sec:lower} reduces $v$ to the normal form $u^2/(1+\mu u)$. It also shows that the lower bound
follows from two results: Theorem~\ref{thmL} (the top invariant lifts), proved in Section~\ref{sec:L}, and
Theorem~\ref{cyc} (cyclicity), proved in Section~\ref{sec:cyc}.
\item Section~\ref{sec:proof} assembles the proof.
\end{itemize}

\section{Reduction to a derivation}\label{sec:reduction}

Throughout, $h_k$ has weight $k$, $\La^n$ is the weight-$n$ part of $\La$, and $\La_{n,d}=\La^n\cap\Q[h]_{\deg d}$
is the part of length $d$. For a formal power series $V(u)=\sum_{j\ge2}V_ju^j$, let $D_V$ be the derivation of
$\La$ given by
\[
  D_V(h_k)=[t^k]\,V(u(t)),\qquad u(t)=\sum_{k\ge1}h_kt^k .
\]
$D_V$ preserves weight and raises length by at least one. Thus it acts nilpotently on each $\La^n$.

\begin{proposition}\label{R1}
Let $T$ be the algebra automorphism of $\La$ with $T(h_k)=p_k/k$. Then $x\in\ker(R_n-D_n)$ if and only if
$g=\sum_\mu x_\mu h_\mu$ satisfies $T(g)=g$. Moreover $\ker(T-I)=\ker D_v$ on $\La^n$, where $v(u)\partial_u$
is the unique formal vector field with $v=O(u^2)$ whose time-one flow is $u\mapsto\log(1+u)$.
\end{proposition}

\begin{proof}
$R_{\lambda\mu}$ counts the maps from the parts of $\lambda$ to the parts of $\mu$ whose fibre sums are the parts of
$\mu$. For $\mu=\lambda$ such a map permutes parts of equal size, so $R_{\lambda\lambda}=\prod_im_i(\lambda)!$.

Use the Hall inner product, for which $\langle h_\lambda,m_\mu\rangle=\delta_{\lambda\mu}$ and
$\langle p_\lambda,p_\mu\rangle=z_\lambda\delta_{\lambda\mu}$, where $z_\lambda=\prod_ii^{m_i}m_i!$. For
$g=\sum_\mu x_\mu h_\mu$ we have $\langle p_\lambda,g\rangle=\sum_\mu R_{\lambda\mu}x_\mu$ and also
$\langle p_\lambda,g\rangle=z_\lambda[p_\lambda]g$. Hence $(R_n-D_n)x=0$ if and only if
$z_\lambda[p_\lambda]g=R_{\lambda\lambda}[h_\lambda]g$ for all $\lambda$, that is,
$[h_\lambda]g=(\prod_j\lambda_j)[p_\lambda]g$. This says exactly that $g=\sum c_\lambda h_\lambda=\sum c_\lambda
q_\lambda$ with $q_\lambda=\prod_jp_{\lambda_j}/\lambda_j$, i.e.\ $T(g)=g$. All steps are reversible.

Since $\sum_kq_kt^k=\log H(t)$ with $H=1+\sum h_kt^k$, $T$ is induced by the pointwise map $u\mapsto\log(1+u)$.
Thus $T-I$ raises length and is nilpotent on $\La^n$. So $\log T$ is a finite sum and a derivation, and
$\ker\log T=\ker(T-I)$, because $\log T=(T-I)\cdot(\text{a polynomial in }T-I\text{ with constant term }1)$. By the
chain rule $D_v(g(u))=g'(u)v(u)$, $\exp(sD_v)$ acts on $u(t)$ as the time-$s$ flow of $v\partial_u$. So
$\exp(D_v)$ and $T$ agree on every $h_k$, and $D_v=\log T$.
\end{proof}

One computes $v=-\tfrac12u^2+\tfrac1{12}u^3+O(u^4)$. In particular the leading part of $D_v$ is $-\frac12D$, where
\[
  D:=D_{u^2},\qquad D(h_k)=\sum_{i+j=k,\ i,j\ge1}h_ih_j .
\]
$D$ is homogeneous: it maps $\La_{n,d}$ to $\La_{n,d+1}$. Write $K_{n,d}=\ker(D\colon\La_{n,d}\to\La_{n,d+1})$ and
$K_d=\bigoplus_nK_{n,d}$.

\section{The kernel of the leading derivation}\label{sec:thmC}

\begin{theorem}\label{thmC}
For every $d\ge0$,
\[
  \sum_n\dim K_{n,d}\;q^n=f^{(d,d)}(q)=q^d\,\qbin{2d}{d}\frac{1}{[d+1]_q},
\]
where $f^{(d,d)}(q)=\sum_{T\in\mathrm{SYT}(d,d)}q^{\mathrm{maj}(T)}$ is the fake degree of the Specht module
$S^{(d,d)}$, i.e.\ the graded multiplicity of $S^{(d,d)}$ in the coinvariant algebra
$R_{2d}=\Q[x_1,\dots,x_{2d}]/(\La^+_{2d})$.
\end{theorem}

\begin{proof}
Fix $n$ and let $m\ge\max(n+1,3)$. Put $V=\Q[t]/(t^m)$, and give $t^i$ (and the $i$-th coordinates below) weight
$i$. All statements in weight $n$ involve only coefficients of index at most $n$, so they do not depend on $m$.

\emph{(a) Pairs.}
\begin{itemize}
\item For $f\in\La$ of length $d$, put $F(v,w)=f(v/w)$ for $v\in tV$ and $w\in1+tV$. Then $F$ is polynomial in the
coefficients of $v$ and $w$, homogeneous of degree $d$ in $v$, and invariant under $(v,w)\mapsto(rv,rw)$ for $r$
in the unit group $U_1=1+tV$.
\item Since $v/(w-\varepsilon v)=u/(1-\varepsilon u)$ and $u\mapsto u/(1-\varepsilon u)$ is the flow of
$u^2\partial_u$, we have $Df=0$ if and only if $F$ is invariant under the unipotent map
$(v,w)\mapsto(v,w-\varepsilon v)$.
\item Conversely, $f(u)=F(u,1)$.
\end{itemize}

\emph{(b) Pl\"ucker extension.} Let $F$ be unipotent-invariant, and let $\omega=a\wedge b\in\wedge^2V$ be
decomposable with some $\omega_{0j}\ne0$.
\begin{itemize}
\item If $a_0=b_0=0$ then every $\omega_{0j}=0$. So the plane $\langle a,b\rangle$ contains $w$ with $w_0=1$
and $v$ with $v_0=0$, scaled so that $\omega=v\wedge w$. The pair $(v,w)$ is unique up to $w\mapsto w+\beta v$,
so $\Phi(\omega):=F(v,w)$ is well defined and homogeneous of degree $d$.
\item On $\{\omega_{0i}\ne0\}$ we have $v_j=-\omega_{0j}$, and normalizing $w_i=0$ gives
$w_j=-\omega_{ij}/\omega_{0i}$. So $\Phi\in A[1/\omega_{0i}]$ for every $i$, where $A$ is the homogeneous
coordinate ring of the cone over $\Gr(2,V)$. The local expressions agree on dense open sets, so they define one
element of $\operatorname{Frac}(A)$.
\item $A$ is the ring of $SL_2$-invariants of $\Q[V\oplus V]$, hence normal. The common zero locus of the
$\omega_{0i}$ on the cone is the cone over $\Gr(2,tV)$, of dimension $2m-5$ inside a cone of dimension $2m-3$. So
no height-one prime contains all the $\omega_{0i}$, and $\bigcap_iA[1/\omega_{0i}]=A$. Hence $\Phi\in A_d=:\Pl_d$.
\item Invariance of $F$ under the connected unipotent group $U_1$ is equivalent to $T_k\Phi=0$ for all $k\ge1$,
where $T_k$ is the derivation induced by multiplication by $t^k$ on $V$. Conversely every
$\Phi\in\Pl_d^T:=\{\Phi:T_k\Phi=0\ \forall k\}$ gives such an $F$, by $F(v,w)=\Phi(v\wedge w)$.
\end{itemize}
The substitution $t\mapsto\lambda t$ scales $\omega_{ij}$ by $\lambda^{i+j}$ and $h_k$ by $\lambda^k$. So these
maps preserve weight, and $K_{n,d}\cong(\Pl_d^T)_n$.

\emph{(c) Schur--Weyl.} $\Pl_d\cong S_{(d,d)}(V)^*$ \cite{Ful}. Set
$M=\Hom_{S_{2d}}(S^{(d,d)},V^{\otimes2d})\cong S_{(d,d)}(V)$, a $GL(V)$-equivariant isomorphism. Here
$V^{\otimes2d}=\Q[x_1,\dots,x_{2d}]/(x_i^m)$, graded by total degree. Multiplication by $t^k$ on $V$ acts on
$V^{\otimes2d}$ as multiplication by the power sum $p_k(x)$. Weight-$n$ functions pair with weight-$n$ vectors.
Hence $(\Pl_d^T)_n\cong((M/\La^+M)_n)^*$, where $\La^+$ is the ideal generated by the $p_k$, and in degrees
below $m$ the truncation is invisible.

\emph{(d) Chevalley.} $\Q[x_1,\dots,x_{2d}]\cong\La_{2d}\otimes R_{2d}$ as graded $S_{2d}$- and $\La_{2d}$-modules
\cite{Che}. So $M/\La^+M\cong\Hom_{S_{2d}}(S^{(d,d)},R_{2d})$. Its graded dimension is the fake degree
$f^{(d,d)}(q)=q^{b(\lambda)}[2d]_q!/\prod_{c}[h(c)]_q$ \cite{Sta79}. For $\lambda=(d,d)$ we have $b=d$ and hook
product $[d+1]_q![d]_q!$, which gives the stated formula.
\end{proof}

\section{Dyck paths}\label{sec:dyck}

Write a Dyck path of semilength $d$ as $U^{a_1}D^{b_1}\cdots U^{a_k}D^{b_k}$, with $A_i=a_1+\dots+a_i$,
$B_i=b_1+\dots+b_i$ and $B_0=0$. Its peaks lie at $x_i=A_i+B_{i-1}$ ($1\le i\le k$), and its valleys at
$v_i=A_i+B_i$ ($1\le i\le k-1$). Hence $\sum_{i=1}^kx_i-\sum_{i=1}^{k-1}v_i=A_k+B_0=d$.

The sum of valley positions is the major index of the path with descents at $DU$, so MacMahon's theorem
\cite{MacM} gives $\sum_{P}q^{\sum v_i}=C_d(q)$. Therefore
\[
  \sum_{P\text{ of semilength }d}q^{\text{peak sum}}=q^dC_d(q),\qquad\text{so}\qquad
  a(n)=\sum_d[q^{n-d}]\,C_d(q)=\sum_d\dim K_{n,d}
\]
by Theorem~\ref{thmC}.

\section{Strategy for the full derivation}\label{sec:lower}

\subsection{Normal form}

\begin{lemma}\label{normal}
Let $v=au^2+bu^3+O(u^4)$ with $a\ne0$, and put $\mu=-b/a$ and $V_\mu=u^2/(1+\mu u)$. There is a formal
$\psi(u)=u+\sum_{j\ge2}c_ju^j$ with $\psi'(u)\,v(u)=a\,V_\mu(\psi(u))$. If $\Theta$ is the weight-preserving
automorphism $h_k\mapsto[t^k]\psi(u(t))$ of $\La$, then $\Theta^{-1}D_v\Theta=aD_{V_\mu}$.
\end{lemma}

\begin{proof}
Compare coefficients of $u^{j+1}$. The unknown $c_j$ enters the left side as $ja\,c_j$ and the right side as
$2a\,c_j$, plus terms in $c_i$ with $i<j$.
\begin{itemize}
\item For $j=2$ the $c_2$ terms cancel, and the $u^3$ equation is $b=-a\mu$, which holds by the choice of $\mu$.
\item For $j\ge3$ the coefficient of $c_j$ is $a(j-2)\ne0$, so $c_j$ is determined recursively.
\end{itemize}
$D_v\circ\Theta$ and $\Theta\circ aD_{V_\mu}$ are $\Theta$-twisted derivations that agree on $u(t)$, so they are
equal. (Compare the formal normal form of a parabolic vector field \cite{Tak}.)
\end{proof}

For the vector field $v$ of Proposition~\ref{R1}, $a=-\frac12$ and $\mu=\frac16$. Hence
$\dim\ker D_v=\dim\ker D_{V_{1/6}}$ on each $\La^n$. Since $\Theta(h_k)=h_k+(\text{longer})$, $\Theta$ preserves
each $F_{\ge d}$ and induces the identity on $F_{\ge d}/F_{\ge d+1}$.

\subsection{Initial forms and source operators}

For $k\ge1$ let $L_k$ be the derivation of $\La$ with $L_k(h_m)=(m-k)h_{m-k}$, where $h_j=0$ for $j\le0$. $L_k$
is induced by the reparametrization $t\mapsto t+\varepsilon t^{k+1}$ of the series variable, since
$u(t+\varepsilon t^{k+1})=u+\varepsilon t^{k+1}u'$. It preserves length and lowers weight by $k$.
Reparametrizations of $t$ commute with every pointwise map $u\mapsto\phi(u)$, so each $L_k$ commutes with every
$D_V$. Since $[L_a,L_b]=(a-b)L_{a+b}$, the operators $L_1,L_2$ generate all $L_k$. Let $U(S)$ denote the algebra
they generate.

\begin{proposition}\label{initial}
Fix $\mu$, let $\pi_d\colon F_{\ge d}\to\La_{\cdot,d}$ be the projection to length $d$, and put
\[
  I_{n,d}=\pi_d\bigl(\ker D_{V_\mu}\cap F_{\ge d}\cap\La^n\bigr),\qquad I_d=\textstyle\bigoplus_nI_{n,d}.
\]
Then $\dim(\ker D_{V_\mu}\cap\La^n)=\sum_d\dim I_{n,d}$, $I_{n,d}\subseteq K_{n,d}$, and $I_d$ is a $U(S)$-stable
subspace of $K_d$.
\end{proposition}

\begin{proof}
The kernel of $\pi_d$ on $\ker D_{V_\mu}\cap F_{\ge d}$ is $\ker D_{V_\mu}\cap F_{\ge d+1}$, which gives the
dimension count. Since $D_{V_\mu}=D-\mu D_{u^3}+\dots$ and $D_{u^j}$ raises length by $j-1$, the length-$(d+1)$
part of $D_{V_\mu}f$ is $D(\pi_df)$, so $I_{n,d}\subseteq K_{n,d}$. Each $L_k$ preserves $F_{\ge d}$ and length
and commutes with $D_{V_\mu}$, so $I_d$ is $U(S)$-stable.
\end{proof}

By Theorem~\ref{thmC}, $K_{n,d}=0$ for $n>d^2$, and $K_{d^2,d}$ is one-dimensional. Let $T_d$ span it.

\begin{theorem}[Lifting the top invariant]\label{thmL}
For every $d\ge1$ and every $\mu$, $T_d\in I_{d^2,d}$. That is, some element of $\ker D_{V_\mu}\cap\La^{d^2}$ has
no component of length $<d$ and has length-$d$ component $T_d$.
\end{theorem}

\begin{theorem}[Cyclicity]\label{cyc}
For every $d\ge1$, $K_d=U(S)\cdot T_d$.
\end{theorem}

These are proved in Sections~\ref{sec:L} and~\ref{sec:cyc}.

\section{Proof of Theorem~\ref{thmL}}\label{sec:L}

\subsection{Invariant coordinates}

Let $W(u)=-1/u+\mu\log u$. Then $W'(u)V_\mu(u)=1$, so $D_{V_\mu}$ sends the series $W(u(t))$ to $1$. Write
\[
  c'_k=[t^k]\,\frac{h_1t}{u(t)},\qquad l_k=[t^k]\log\frac{u(t)}{h_1t},\qquad e_k=c'_{k+1}-\mu h_1l_k\quad(k\ge1).
\]
Then $\La[1/h_1]=\Q[h_1^{\pm1},c'_1,e_1,e_2,\dots]$, and in these coordinates $D_{V_\mu}=-h_1\,\partial/\partial c'_1$.
So $\ker D_{V_\mu}$ on $\La[1/h_1]$ is $\Q[h_1^{\pm1},e_1,e_2,\dots]$, and $e_k$ has weight $k+1$. Every element
of $\ker D_{V_\mu}\cap\La^n$ is $\sum_ah_1^aG_a(e)$ with $G_a\in\Q[x_1,x_2,\dots]$ of weight $n-a$, where
$\operatorname{wt}x_k=k+1$.

\subsection{A factorization model}

Treat $\mu$ as an indeterminate. Let $\mathcal R=\Q[\mu][h_1,\zeta,A_1,A_2,\dots]$, and embed $\La[\mu]$ by
writing $u(t)/t=(h_1-\zeta t)/A(t)$ with $A=1+A_1t+A_2t^2+\cdots$.

\begin{lemma}\label{fact}
$\La[\mu]=\La[\mu][1/h_1]\cap\mathcal R$.
\end{lemma}

\begin{proof}
Modulo $h_1$ the embedding sends $h_k\mapsto-\zeta[t^{k-2}]A^{-1}$, which is injective on $\Q[h_2,h_3,\dots]$. So if
$h_1$ divides the image of $F$ in $\mathcal R$, then $h_1$ divides $F$ in $\La$. Induct on the power of $h_1$.
\end{proof}

In $\mathcal R$ we have $h_1t/u=A\sum_j(\zeta/h_1)^jt^j$ and $\log(u/h_1t)=\log(1-\zeta t/h_1)-\log A$. This gives
\[
  h_1^{k+1}e_k=\Gamma_{k+1}+\mu h_1^2\Bigl(\frac{\zeta^k}{k}+h_1^k\lambda_k\Bigr),\qquad
  \Gamma_k=\sum_{j\le k}A_jh_1^j\zeta^{k-j},\quad \lambda_k=[t^k]\log A .
\]
Put $d_j=A_jh_1^j$ and $\nu=\mu h_1^2$, and set $\zeta=1$. Then $h_1^{k+1}e_k=S_{k+1}+\nu(1/k+\pi_k)$, where
$S_k=1+d_1+\dots+d_k$ and $\pi_k=[s^k]\log(1+\sum_jd_js^j)$.

Setting $\zeta=1$ loses nothing. With $\operatorname{wt}\zeta=2$, $\operatorname{wt}d_j=2j$ and
$\operatorname{wt}\nu=2$, each $h_1^{k+1}e_k$ is homogeneous, so the $\zeta$-exponent of every monomial is fixed by
its $(d,\nu)$-part. Moreover distinct monomials in $(d,\nu)$ map to distinct monomials in $(A,\mu,h_1)$. So the
$h_1$-adic valuation of a polynomial $P(d,\nu)$ equals its \emph{order}: the least weight of its monomials, with
$\operatorname{wt}d_j=j$ and $\operatorname{wt}\nu=2$. Combining with Lemma~\ref{fact}:

\begin{lemma}[d-model]\label{dmodel}
Let $G_{a_0+i}\in\Q[x]$ have weight $n-a_0-i$. Then
$\sum_i\mu^ih_1^{a_0+i}G_{a_0+i}(e)\in\La[\mu]^n$ if and only if
$\sum_i\nu^iG_{a_0+i}\bigl(S_{k+1}+\nu(1/k+\pi_k)\bigr)$ has order at least $n-2a_0$. (Here $x_k$ is replaced by
$S_{k+1}+\nu(1/k+\pi_k)$.)
\end{lemma}

Let $z_1=x_1-1$ and $z_k=x_k-x_{k-1}$ for $k\ge2$, with weights $1,3,4,5,\dots$. Under $x_k\mapsto S_{k+1}$ we
have $z_1\mapsto d_1+d_2$ and $z_k\mapsto d_{k+1}$, so the order of $Q(S)$ equals the weighted order of $Q$ at
$x=(1,1,\dots)$.

\begin{lemma}\label{harmless}
For every polynomial $Q(x,\nu)$, $\operatorname{ord}Q(S+\nu\pi)\ge\operatorname{ord}Q(S)$.
\end{lemma}

\begin{proof}
Expand by Taylor's formula. $\partial/\partial x_k=\partial_{z_k}-\partial_{z_{k+1}}$ lowers weighted order by at
most $k+2$, and $\nu\pi_k$ has order at least $k+2$.
\end{proof}

\subsection{The forest--Heine polynomial}

Let $m=d-1$, let $\rho$ be the moment functional $\langle\rho,s^k\rangle=x_{k+1}$, and let
$\Delta=\prod_{i<j}(s_i-s_j)$. By Heine's formula, $G_d:=\det(x_{i+j-1})_{i,j\le m}=\frac1{m!}\langle\rho^{\otimes
m},\Delta^2\rangle$.

\begin{definition}
$F(x,\nu)=\frac1{m!}\sum_{\mathcal F}(-\nu)^{\#\mathrm{children}}\int\Delta^2$, where $\mathcal F$ runs over the
rooted forests on $[m]$. Each root $i$ carries $\rho(ds_i)$. Each child $c$ with parent $p$ is integrated over
$[0,s_p]$ with Lebesgue measure, using the polynomial rule $s^k\mapsto s_p^{k+1}/(k+1)$.
\end{definition}

The $\nu^i$ term of $F$ has weight $m^2+m-i$, and $F(x,0)=G_d$. Put
\[
  F_d:=\sum_i\mu^ih_1^{d+i}\,[\nu^i]F(e).
\]

\begin{proposition}\label{Fprop}
$F(S+\nu(1/k+\pi_k),\nu)$ has order at least $m^2-1$. Consequently $F_d\in\La[\mu]^{d^2}\cap\ker D_{V_\mu}$.
\end{proposition}

\begin{proof}
\emph{Step 1.} Apply Lemma~\ref{harmless} to $Q(x,\nu)=F(x+\nu v,\nu)$ with $v_k=1/k$. It suffices to bound the
order of $F(S+\nu v,\nu)$. Since $\langle\mathrm{Leb}[0,1],s^k\rangle=1/(k+1)$, the vector $S+\nu v$ is the moment
sequence of $\rho+\nu\,\mathrm{Leb}[0,1]$.

\emph{Step 2.} Write $\langle\rho,s^k\rangle=S_{k+2}=B-\gamma_k$, where $B=1+\sum_jd_j$ has order $0$ and
$\gamma_k=\sum_{i\ge k+3}d_i$ has order at least $k+3$. It suffices to treat $B$, $\gamma_k$ and $\nu$ as
independent variables of orders $0$, $k+3$ and $2$, and to show that every surviving monomial has order at least
$m^2-1$.

\emph{Step 3.} Realize this with a parameter $\tau$, where $0<\tau<1/C$.
\begin{itemize}
\item Take $\rho=B\delta_1-J_\tau$, with $J_\tau(g)=\tau^3\int g(\tau y)\,\omega(dy)$, where $\omega$ is a
polynomial density on $[0,C]$ with the finitely many prescribed moments $\gamma_k$. Replace $\nu$ by $\tau^2\nu$.
\item Then $E(\tau):=F(\text{moments of }\rho+\tau^2\nu\,\mathrm{Leb})$ is the polynomial evaluated at
$(B,\tau^{k+3}\gamma_k,\tau^2\nu)$. So it suffices to show $E(\tau)=O(\tau^{m^2-1})$ for all real data.
\end{itemize}

\emph{Step 4.} For a given child set, sum over parent maps. A cycle in the parent map would force
$s_{c_1}<\dots<s_{c_1}$, and ties have measure zero. So the sum of the indicator products is
$\prod_c\mathbf 1_{s_c\ge0}N_c(s)$, where $N_c=\#\{q\ne c:s_q>s_c\}$. Expand $\rho+\nu\,\mathrm{Leb}$ at the roots
and group the Lebesgue roots with the children. This gives
\[
  E=\frac1{m!}\sum_{U\subseteq[m]}\nu^{|U|}\int\prod_{i\notin U}\rho(ds_i)\prod_{i\in U}
  \bigl(\mathbf 1_{[0,1]}(s_i)-N_i(s)\bigr)\mathbf 1_{s_i\ge0}\,ds_i\ \Delta(s)^2 .
\]

\emph{Step 5.} The points are of three kinds: $e$ points at $1$ (from $B\delta_1$); cluster points (roots from
$J_\tau$, and $U$-points in $[0,C\tau]$); and $h$ high $U$-points in $(C\tau,1]$.
\begin{itemize}
\item If $e\ge2$ then $\Delta=0$.
\item A $U$-point above $1$ that is maximal among the $U$-points has density $0$, so every $U$-point lies in
$[0,1]$.
\item For a high $U$-point, $N_i=e+\#\{\text{high }U\text{-points above it}\}$. So the topmost high point has
density $1-e$ and the next one has density $-e$. Hence $e=1$ forces $h=0$, and $e=0$ forces $h\le1$.
\end{itemize}

\emph{Step 6.} Each cluster point costs $\tau^3$, each high point costs $\tau^2$, and each pair of cluster points
contributes $(s_i-s_j)^2=O(\tau^2)$. With $e+c+h=m$, a configuration is $O(\tau^{3c+c(c-1)+2h})$. The three
surviving cases $(e,h)=(1,0),(0,1),(0,0)$ give the exponents $m^2-1$, $m^2+1$ and $m^2+2m$.

\emph{Conclusion.} By Lemma~\ref{dmodel} with $a_0=d$ and $n=d^2$, the order bound $m^2-1=n-2a_0$ shows that
$F_d\in\La[\mu]$. Finally $F_d\in\ker D_{V_\mu}$, because it is a polynomial in $h_1$ and the $e_k$.
\end{proof}

\begin{proof}[Proof of Theorem~\ref{thmL}]
Let $\sigma_\lambda$ be the automorphism $h_k\mapsto\lambda h_k$. Then $c'_k$ and $l_k$ are $\sigma$-invariant, and
$\sigma_\lambda e_k(\mu)=e_k(\lambda\mu)$, so $\sigma_\lambda F_d(\mu)=\lambda^dF_d(\lambda\mu)$. Writing
$F_d(\mu)=\sum_j\mu^j\Phi_j$, each $\Phi_j$ therefore has length exactly $d+j$. So $F_d$ has no component of length
$<d$, and its length-$d$ component is $F_d(0)=h_1^dG_d(e|_{\mu=0})$. This lies in $K_{d^2,d}$: it is a nonzero
element, by the case $\nu=0$ of Proposition~\ref{Fprop}, and $K_{d^2,d}$ is one-dimensional. Specializing $\mu$ to
any value gives an element of $\ker D_{V_\mu}\cap\La^{d^2}$ with initial form $T_d$, up to scalar.
\end{proof}

\section{Proof of Theorem~\ref{cyc}}\label{sec:cyc}

\subsection{Duality with the coinvariant side}

On $\Q[x_1,\dots,x_N]$ let $E_k=\sum_ix_i^{k+1}\partial/\partial x_i$ for $k\ge1$. $E_k$ is a derivation, commutes
with $S_N$, and maps $p_r$ to $rp_{r+k}$. So it preserves $\La^+\Q[x]$ and acts on $R_N$ and on
$N_\lambda:=\Hom_{S_N}(S^\lambda,R_N)$.

\begin{lemma}\label{dual}
Under the isomorphism $K_{n,d}\cong(N_{(d,d)})_n^*$ of Theorem~\ref{thmC}, $L_k$ corresponds to $\pm E_k^{\,T}$.
\end{lemma}

\begin{proof}
Follow the proof of Theorem~\ref{thmC}, with $m$ large enough to serve all weights at once.
\begin{itemize}
\item Since $(v\circ\phi)/(w\circ\phi)=(v/w)\circ\phi$, $L_k$ corresponds to the derivation of pair functions
induced by $\xi_k=t^{k+1}\partial_t\in\mathfrak{gl}(V)$, acting on both $v$ and $w$. This preserves $\{v_0=0\}$ and
$\{w_0=1\}$.
\item Since $[\xi_k,t^j\cdot]=jt^{j+k}\cdot$, $\xi_k$ normalizes the Lie algebra of $U_1$, so it preserves
$\Pl_d^T$.
\item The isomorphism $\Pl_d\cong M^*$ is $\mathfrak{gl}(V)$-equivariant, and $\xi_k$ acts on
$V^{\otimes2d}=\Q[x]/(x_i^m)$ as $E_k$, because $\xi_kt^a=at^{a+k}$.
\item $E_k$ preserves $\La^+M$ and commutes with $S_{2d}$, and the identification
$M/\La^+M\cong N_{(d,d)}$ is $E$-equivariant.
\end{itemize}
The sign does not affect spans.
\end{proof}

\begin{lemma}\label{socle}
Let $N$ be a finite-dimensional graded module on which each $E_k$ raises degree by $k$. Suppose its top degree is
one-dimensional, and let $\theta$ span $(N^*)_{\mathrm{top}}$. Then $U(E^T)\cdot\theta=N^*$ if and only if
$\soc(N):=\{x:E_kx=0\ \forall k\}$ lies in the top degree.
\end{lemma}

\begin{proof}
$(U\cdot\theta)^\perp$ is a graded $E$-submodule meeting the top degree in $0$. A nonzero graded $E$-submodule
contains a nonzero homogeneous socle vector, because $E$ acts nilpotently. Conversely, a socle vector of lower
degree spans a submodule inside $(U\cdot\theta)^\perp$.
\end{proof}

The top weight of $K_d$ is $d^2$, the top degree of $f^{(d,d)}(q)$, with coefficient $1$. By
Lemmas~\ref{dual} and~\ref{socle}, Theorem~\ref{cyc} is equivalent to: \emph{$\soc_E(N_{(d,d)})$ is its top degree
$d^2$.}

\subsection{Gorenstein duality}

Let $N=2d$ and $\mathrm{top}=\binom{N}{2}$. Define $c\colon\Q[x]_{\mathrm{top}}\to\Q$ by $\mathcal A(g)=c(g)\Delta$,
where $\mathcal A$ is antisymmetrization and $\Delta$ the Vandermonde. Then $c$ vanishes on $\La^+\Q[x]$, and
$\langle f,g\rangle=c(fg)$ is a perfect pairing on $R_N$ that pairs $\lambda$-isotypic parts with
$\lambda'$-isotypic parts. Since $\mathcal A(E_kh)=E_k(\mathcal Ah)=0$ for $\deg h<\mathrm{top}$, we get
$c(E_kh)=0$, so $\langle E_kf,g\rangle=-\langle f,E_kg\rangle$. Hence $N_{(d,d)}$ and $N_{(2^d)}$ are perfectly
paired, between degrees $j$ and $\mathrm{top}-j$, with $E$ skew. Then
$\soc(N_{(d,d)})^\perp=\sum_kE_kN_{(2^d)}$, and graded Nakayama gives:

\begin{corollary}
Theorem~\ref{cyc} holds if and only if $N_{(2^d)}$ is generated under $E$ by its bottom degree $d(d-1)$. That
degree is spanned by the Specht vector $\Delta(x_1..x_d)\Delta(x_{d+1}..x_{2d})$.
\end{corollary}

\subsection{The Grassmannian model}

Let $X=\{x_1..x_d\}$, $Y=\{x_{d+1}..x_{2d}\}$ and $H=S_d\times S_d$.
\begin{itemize}
\item $\operatorname{Ind}_H^{S_{2d}}\mathrm{sgn}=\bigoplus_{i\le d}S^{(2^i,1^{2d-2i})}$. So the sign-isotypic part
$R^{\mathrm{sgn}_H}$ is $\bigoplus_iN_{(2^i,1^{2d-2i})}$, $E$-equivariantly.
\item Multiplication by $\Delta_X\Delta_Y$ is a bijection
$R^H=(\La_X\otimes\La_Y)/\La^+\to R^{\mathrm{sgn}_H}$. Here $R^H\cong G_{d,2d}:=H^*(\Gr(d,2d))=\La/I_{d,2d}$, where
$I_{j,n}=\operatorname{span}\{s_\lambda:\lambda\not\subseteq j\times(n-j)\}$.
\item One computes $E_k\Delta(Z)=\Delta(Z)\sum_{i<j}h_k(z_i,z_j)$, that is
$E_k\Delta(Z)/\Delta(Z)=(|Z|-\tfrac{k+1}2)p_k(Z)+\tfrac12\sum_{a+b=k}p_a(Z)p_b(Z)$ with $a,b\ge1$. Since
$p_a(Y)\equiv-p_a(X)$ in $R$, we get $E_k(\Delta_X\Delta_Yf)=\Delta_X\Delta_Y\,\fE_kf$ with
\[
  \fE_k:=E_k+Q_k,\qquad Q_k=\sum_{a+b=k,\ a,b\ge1}p_ap_b .
\]
\item $n(\lambda)$ strictly decreases along dominance, so the degree-$d(d-1)$ part of $R^{\mathrm{sgn}_H}$ is the
bottom of $N_{(2^d)}$. It corresponds to $1\in G_{d,2d}$, and $\dim N_{(2^d)}=f^{(2^d)}=\mathrm{Cat}(d)$.
\end{itemize}

\begin{corollary}\label{grassred}
Theorem~\ref{cyc} holds if and only if $\dim U(\fE)\cdot1\ge\mathrm{Cat}(d)$ in $G_{d,2d}$.
\end{corollary}

On $\La=\Q[p_1,p_2,\dots]$, $E_k$ acts by $p_r\mapsto rp_{r+k}$, and $Q_k$ by multiplication. Both preserve each
$I_{j,n}$, because $I_{j,n}$ is the kernel of the $E$-equivariant map $\La\to R_n^{S_j\times S_{n-j}}$.

\subsection{Twisted box modules}

For $c\in\mathbb Z$ put $\fE^{(c)}_k=E_k+Q_k-c\,p_k$, acting on $\La$ and on every $G_{j,n}=\La/I_{j,n}$. Let
$a_c(j,n)=\dim U(\fE^{(c)})\cdot1\subseteq G_{j,n}$ and $L^{(c)}=U(\fE^{(c)})\cdot1\subseteq\La$. Write
$f^{(a,b)}$ for the number of standard tableaux of shape $(a,b)$, with $f^{(a,b)}=0$ if $b>a$.

\begin{lemma}\label{exact}
Let $m=n-j$.
\begin{enumerate}
\item[(i)] The quotient $\rho\colon G_{j,n}\to G_{j,n-1}$ commutes with every $\fE^{(c)}$. Multiplication by
$h_m$ induces an injection $\iota\colon G_{j-1,n-1}\to G_{j,n}$, $s_\mu\mapsto s_{(m,\mu)}$, with image
$\ker\rho$, and $\fE^{(c)}_k\iota=\iota\fE^{(c+1)}_k$.
\item[(ii)] The quotient $\rho'\colon G_{j,n}\to G_{j-1,n-1}$ commutes with $\fE^{(c)}$. Multiplication by $e_j$
induces an injection $\iota'\colon G_{j,n-1}\to G_{j,n}$, $s_\mu\mapsto s_{\mu+(1^j)}$, with image $\ker\rho'$, and
$\fE^{(c)}_k\iota'=\iota'\fE^{(c-1)}_k$.
\end{enumerate}
\end{lemma}

\begin{proof}
\begin{itemize}
\item The statements about images follow from Pieri's rule.
\item For (i): $E_kH(z)=H(z)\sum_rp_{r+k}z^r$ gives $E_kh_m=\sum_{s>k}p_sh_{m+k-s}$. By Newton's identity,
$E_kh_m+p_kh_m=(m+k)h_{m+k}-\sum_{s<k}p_sh_{m+k-s}$. Every $h$ on the right has index $>m$, so it lies in
$I_{j,n}$. Hence $\fE^{(c)}_k(h_mf)=h_m(\fE^{(c)}_k-p_k)f$ in $G_{j,n}$.
\item For (ii): similarly $E_ke_j-p_ke_j=(-1)^k[(j+k)e_{j+k}-\sum_{s<k}(-1)^{s-1}p_se_{j+k-s}]$. Every $e$ here
has index $>j$, so it lies in $I_{j,n}$.
\end{itemize}
\end{proof}

\begin{lemma}\label{members}
\begin{itemize}
\item[(a)] $e_j=\frac{(-1)^j}{j!}(\fE^{(1)}_1)^j\cdot1\in L^{(1)}$.
\item[(b)] $C_m:=\frac12(h_m+(-1)^me_m)\in L^{(0)}$ for every $m$ (Theorem~\ref{claimA}). Hence
$h_m\in L^{(0)}+I_{j,j+m}$ whenever $1\le j<m$, because $e_m\in I_{j,j+m}$.
\end{itemize}
\end{lemma}

\begin{proof}[Proof of (a)]
With $\phi(z)=\sum p_rz^r/r$ we have $\fE^{(c)}_1e^{\lambda\phi(z)}=\partial_ze^{\lambda\phi(z)}-(\lambda+c)p_1
e^{\lambda\phi(z)}$, since $Q_1=0$. Take $\lambda=-1$ and $c=1$. Since $e^{-\phi(z)}=\sum_r(-1)^re_rz^r$, we get
$\fE^{(1)}_1\varepsilon_{r-1}=r\varepsilon_r$ for $\varepsilon_r=(-1)^re_r$.
\end{proof}

Moreover $\omega I_{j,n}=I_{n-j,n}$, and $\omega\fE^{(0)}_k\omega=(-1)^k\fE^{(0)}_k$, so $a_0(j,n)=a_0(n-j,n)$.

\begin{theorem}\label{induction}
$a_0(j,n)\ge f^{(n-j,j)}$ for $n\ge2j$, and $a_1(j,n)\ge f^{(n-j,j)}$ for $n\ge2j+1$.
\end{theorem}

\begin{proof}
Induct on $n$; for $j=0$, $G_{0,n}=\Q$. For a surjection with kernel $\iota(G')$ as in Lemma~\ref{exact},
$\dim U\cdot1\ge\dim\rho(U\cdot1)+\dim(U\cdot1\cap\ker\rho)$, and $\rho(U\cdot1)$ is the cyclic module in the target.
\begin{itemize}
\item \emph{$c=0$, $2j<n$.} By Lemma~\ref{members}(b), $h_m\in U\cdot1$. Hence $U\cdot1\cap\ker\rho\supseteq
\iota(U^{(1)}\cdot1)$, so $a_0(j,n)\ge a_0(j,n-1)+a_1(j-1,n-1)\ge f^{(n-j-1,j)}+f^{(n-j,j-1)}=f^{(n-j,j)}$.
\item \emph{$c=0$, $n=2j$.} By $\omega$-symmetry and the previous case,
$a_0(j,2j)\ge a_0(j,2j-1)=a_0(j-1,2j-1)\ge f^{(j,j-1)}=f^{(j,j)}$.
\item \emph{$c=1$, $n\ge2j+1$.} $e_j\in U^{(1)}\cdot1$ by Lemma~\ref{members}(a), and Lemma~\ref{exact}(ii) applies.
So $a_1(j,n)\ge a_1(j-1,n-1)+a_0(j,n-1)\ge f^{(n-j,j-1)}+f^{(n-j-1,j)}=f^{(n-j,j)}$.
\end{itemize}
\end{proof}

\subsection{A screening operator}

Let $\phi(z)=\sum_rp_rz^r/r$, $\Phi(z)=z\phi'(z)$ and $H(z)=e^{\phi(z)}$. Let $f[X-w]$ denote the substitution
$p_r\mapsto p_r-w^r$. Put $L=L^{(0)}=U(\fE)\cdot1$ with $\fE_k=E_k+Q_k$.

\begin{lemma}\label{lowerdim}
$\dim L_m\ge p(m)-p(m-1)$.
\end{lemma}

\begin{proof}
$E_k$ preserves $p$-length, and $Q_k$ raises it by exactly $2$. So the top-length part of
$\fE_{\lambda_1}\cdots\fE_{\lambda_l}\cdot1$ is $\prod Q_{\lambda_i}$. For $k\ge2$ the $Q_k$ are algebraically
independent: $\partial Q_k/\partial p_{k-1}=2p_1$, so the Jacobian is triangular. Hence the $\fE_\lambda\cdot1$
with $\lambda\vdash m$ and all parts $\ge2$ are independent, and there are $p(m)-p(m-1)$ of them.
\end{proof}

Define $J\colon\La_m\to\La_{m-1}$ by $Jf=\Res_zH(z)^2f[X-1/z]$.

\begin{lemma}\label{intertw}
$J(\fE_kf)=(\fE_k+2p_k)Jf$ and $J(1)=0$. Hence $J(L)=0$.
\end{lemma}

\begin{proof}
Write $T_w$ for $f\mapsto f[X-w]$ and $D_w=\sum_rrw^{r-1}\partial_{p_r}$.
\begin{itemize}
\item One has $T_wE_kT_w^{-1}=E_k-w^{k+1}D_w$, $D_wT_wf=-\partial_w(T_wf)$, and
$T_wQ_k=Q_k-2\sum_{a+b=k}w^ap_b+(k-1)w^k$.
\item Put $w=1/z$ and $G=H(z)^2T_{1/z}f$. Using $E_k(H^2)=2H^2z^{-k}(\Phi-\sum_{s\le k}p_sz^s)$, these give
\[
  H(z)^2T_{1/z}(\fE_kf)=(\fE_k+2p_k)G+(k-1)z^{-k}G-z^{1-k}\partial_zG .
\]
\item The last two terms equal $-\partial_z(z^{1-k}G)$, which has zero residue.
\item $J(1)=\Res H(z)^2=0$, because $H(z)^2$ is a power series in $z$.
\end{itemize}
\end{proof}

\begin{lemma}\label{Jkill}
$J(C_m)=0$.
\end{lemma}

\begin{proof}
$h_m[X-w]=h_m-wh_{m-1}$, so $Jh_m=-h_{m-1}$. Next $e_m[X-w]=\sum_i(-w)^ie_{m-i}$, so
$\sum_mu^mJe_m=-uE(u)H(-u)^2=-uH(-u)$. This gives $J((-1)^me_m)=h_{m-1}$.
\end{proof}

\begin{lemma}\label{surj}
$J\colon\La_m\to\La_{m-1}$ is surjective.
\end{lemma}

\begin{proof}
Equip $\La$ with the positive definite form $\langle p_\lambda,p_\mu\rangle=\delta_{\lambda\mu}z_\lambda
2^{-\ell(\lambda)}$, for which $p_r^\dagger=\frac r2\partial_{p_r}$. Define
\begin{align*}
  J^\vee f&=\Res_w\,w^{-2}H(w)^{-2}f[X+1/w],\\
  Af&=\Res_z\,z^2H(z)^2f[X-1/z],\\
  A^\vee f&=\Res_w\,w^{-4}H(w)^{-2}f[X+1/w].
\end{align*}
Write $H(z)^2=\exp(\sum2p_rz^r/r)$ and $f[X-1/z]=\exp(-\sum z^{-r}\partial_{p_r})f$, and substitute $z=1/w$. One
finds $J^\dagger=J^\vee$ and $A^\dagger=A^\vee$.

The compositions $JJ^\vee$ and $A^\vee A$ come from the same kernel,
\[
  z^2w^{-2}(z-w)^{-2}H(z)^2H(w)^{-2}f[X+1/w-1/z],
\]
expanded in $|w|<|z|$ and in $|z|<|w|$ respectively. By the formal identity
$\iota_{z,w}(z-w)^{-2}-\iota_{w,z}(z-w)^{-2}=\partial_w\delta(z,w)$ and integration by parts,
\[
  JJ^\vee-A^\vee A=\Res_w\,w^{-2}\partial_z\bigl(z^2H(z)^2H(w)^{-2}f[X+1/w-1/z]\bigr)\big|_{z=w}=2f .
\]
The last equality holds because the only term with nonzero residue is $2w^{-1}f$. Hence
$JJ^\dagger=2+A^\dagger A$, which is positive definite on each $\La_{m-1}$. So $J$ is surjective.
\end{proof}

\begin{theorem}\label{claimA}
$L_m=\ker(J|_{\La_m})$, and $C_m\in L$ for all $m$.
\end{theorem}

\begin{proof}
By Lemma~\ref{surj}, $\dim\ker J|_{\La_m}=p(m)-p(m-1)$. By Lemmas~\ref{lowerdim} and~\ref{intertw}, $L_m\subseteq\ker
J$ and $\dim L_m\ge p(m)-p(m-1)$, so the two are equal. Lemma~\ref{Jkill} then gives $C_m\in L$.
\end{proof}

\begin{remark}
With $p_r=b_r/\sqrt2$, the operators $\fE_k$ are the Virasoro modes $L_{-k}$ of the $c=1$ Fock space. $J$, $J^\vee$, $A$ and $A^\vee$ are zero modes of the level-one $\widehat{\mathfrak{sl}}_2$ currents, and the identity in
Lemma~\ref{surj} is the commutator $[J_0^+,J_0^-]=2$ in momentum $\sqrt2$ \cite{KR}. None of this is needed for
the proof.
\end{remark}

\begin{proof}[Proof of Theorem~\ref{cyc}]
Theorem~\ref{induction} gives $a_0(d,2d)\ge f^{(d,d)}=\mathrm{Cat}(d)$. Apply Corollary~\ref{grassred} and the
reductions above.
\end{proof}

\section{Proof of Theorems~\ref{main} and~\ref{refined}}\label{sec:proof}

By Proposition~\ref{R1}, the nullspace of $R_n-D_n$ is identified with $N=\ker D_v\cap\La^n$. By
Lemma~\ref{normal}, $\Theta^{-1}N=\ker D_{V_\mu}\cap\La^n$ with $\mu=\frac16$. Since $\Theta$ preserves the
filtration $F_{\ge d}$ and induces the identity on its graded pieces, it suffices to prove Theorem~\ref{refined}
for $D_{V_\mu}$.

By Proposition~\ref{initial}, the graded pieces have dimensions $\dim I_{n,d}$, and $I_d\subseteq K_d$ is
$U(S)$-stable. By Theorem~\ref{thmL}, $T_d\in I_d$. By Theorem~\ref{cyc}, $K_d=U(S)\cdot T_d\subseteq I_d$. Hence
$I_{n,d}=K_{n,d}$, and Theorem~\ref{thmC} gives $\dim I_{n,d}=[q^{n-d}]C_d(q)$, which is Theorem~\ref{refined}.
Summing over $d$ and using Section~\ref{sec:dyck} gives $\dim\ker(R_n-D_n)=a(n)$.\qed

\section{Remarks}

\begin{remark}
Theorem~\ref{thmC} identifies the invariants of $u\mapsto u/(1-\varepsilon u)$ acting on power series with the
dual of $\Hom_{S_{2d}}(S^{(d,d)},R_{2d})$. Explicitly, for every functional $\nu$ on $R_{2d}$, the polynomial
$f_\nu(u)=\nu\bigl(\prod_{i=1}^d(u(x_i)-u(y_i))\bigr)$ is an invariant. The reason is that the flow multiplies the
product by an $S_{2d}$-symmetric series with constant term $1$, which acts trivially on $R_{2d}$.
\end{remark}

\begin{remark}
The analogue of the statement ``$\soc_E N_\lambda$ is the top degree'' fails for general shapes. For example, for
$\lambda=(3,2,1)$ the socle of $N_\lambda$ contains a vector one degree below the top. Two-row shapes are special.
\end{remark}

\begin{remark}
All identities in Sections~\ref{sec:L} and~\ref{sec:cyc} were also checked by exact computation (see
Reproducibility):
\begin{itemize}
\item the nullity of $R_n-D_n$, computed from the definition, equals $a(n)$ for $n\le17$;
\item $\dim K_{n,d}$ matches Theorem~\ref{thmC} for $n\le17$, $d\le7$;
\item the lift $F_d$ satisfies $D_{V_\mu}F_d=0$ identically in $\mu$ for $d\le7$;
\item Theorem~\ref{cyc} holds directly in $\La$ for $d\le7$;
\item the identities of Lemmas~\ref{intertw}--\ref{surj} hold for $m\le13$.
\end{itemize}
\end{remark}

\subsection*{Reproducibility}
{\sloppy All code is available at
\url{https://github.com/prajitchennupati/math/tree/main/clifford-stanley-nullity-conjecture}; paths below are
relative to that folder. Every script uses exact arithmetic (python-flint 0.9, SymPy 1.14).
\begin{itemize}
\item \texttt{paper/code/verify.py} computes the nullity of $R_n-D_n$ from the definition, enumerates A129528, checks
MacMahon's identity, and computes $\dim K_{n,d}$.
\item \texttt{paper/code/plucker.py} computes the $T$-invariant Pl\"ucker functions of Theorem~\ref{thmC}.
\item \texttt{paper/code/normal\_form.py} checks Lemma~\ref{normal}.
\item \texttt{paper/code/lift.py} builds $F_d$ from the forest sum and checks Theorem~\ref{thmL}.
\item \texttt{paper/code/cyclicity.py} checks Theorem~\ref{cyc}.
\item \texttt{paper/code/fock.py} checks Lemmas~\ref{exact}--\ref{surj}.
\end{itemize}
This paper has not been refereed.\par}

\begin{thebibliography}{99}
\bibitem{Che} C.~Chevalley, \emph{Invariants of finite groups generated by reflections}, Amer. J. Math. 77 (1955),
778--782.
\bibitem{CS04} P.~Clifford, R.~P.~Stanley, \emph{Bottom Schur functions}, Electron. J. Combin. 11 (2004), \#R67;
arXiv:math/0311382.
\bibitem{Ful} W.~Fulton, \emph{Young Tableaux}, London Math. Soc. Student Texts 35, Cambridge, 1997.
\bibitem{KR} V.~G.~Kac, A.~K.~Raina, \emph{Bombay Lectures on Highest Weight Representations of
Infinite-Dimensional Lie Algebras}, World Scientific, 1987.
\bibitem{Mac} I.~G.~Macdonald, \emph{Symmetric Functions and Hall Polynomials}, 2nd ed., Oxford, 1995.
\bibitem{MacM} P.~A.~MacMahon, \emph{Combinatory Analysis}, Cambridge, 1915--1916.
\bibitem{OEIS} OEIS Foundation, \emph{The On-Line Encyclopedia of Integer Sequences}, sequence A129528,
\url{https://oeis.org/A129528}.
\bibitem{Sta79} R.~P.~Stanley, \emph{Invariants of finite groups and their applications to combinatorics}, Bull.
Amer. Math. Soc. 1 (1979), 475--511.
\bibitem{SymCat} P.~Alexandersson, \emph{The symmetric functions catalog}, \url{https://www.symmetricfunctions.com/schurMisc.htm}
(accessed November 2026).
\bibitem{Tak} F.~Takens, \emph{Normal forms for certain singularities of vectorfields}, Ann. Inst. Fourier 23
(1973), 163--195.
\end{thebibliography}

\end{document}
